\textbf{Data and splits.} scikit-learn digits (\rhval{const/n_images} images, $8{\times}8$, pixel values divided by 16). For each seed $s\in\{0,1,2,3,4\}$ a stratified split gives \rhval{const/n_train} training, \rhval{const/n_val} validation and \rhval{const/n_test} test images; the seed also determines network initialisation, minibatch order, dropout masks and noise. Shift is applied to test images only (and, for the oracle ablation, to the validation images), never to training data.

\textbf{Shift conditions (``tasks'').} Eleven conditions: in-distribution (r0); rotations of 10, 20, 30, 45, 60 degrees (r10--r60); Gaussian noise with $\sigma\in\{\rhval{const/noise_sigma_1:2},\rhval{const/noise_sigma_2:1},\rhval{const/noise_sigma_3:2},\rhval{const/noise_sigma_4:1}\}$ (n0.25--n1.0); and rotation 30 plus noise $\sigma=\rhval{const/noise_sigma_2:1}$ (r30+n0.5). The strongest conditions are not mild: on $8{\times}8$ images accuracy falls to chance level for r45/r60 (Table~\ref{tab:accbrier}), where ECE is therefore essentially the mean confidence.

\textbf{Systems.} All reimplemented in PyTorch in \texttt{method/run.py}: MLP 64--\rhval{const/hidden_units}--\rhval{const/hidden_units}--10 with ReLU; Adam, learning rate \rhval{const/learning_rate:sci0}, batch 64, \rhval{const/epochs} epochs, no weight decay, no early stopping. MLP+TS uses the MLP with $T$ fitted on unshifted validation data; MC dropout trains the same architecture with dropout $p=\rhval{const/dropout_p:1}$ after each hidden layer and averages $S=20$ stochastic passes; the ensemble has $M=5$ MLPs. The single MLP is ensemble member 0. All systems share data, splits, optimiser and epochs, and are scored with the same metric code (15-bin ECE, NLL, Brier, accuracy, mean confidence). Network weights are cached per (architecture, seed) so that the rows of a seed share one trained network; training is deterministic.

\textbf{Tuning budget and tests.} None: every hyperparameter was fixed before the runs and not selected on test data (validation data are used only to fit $T$). The sweeps in Section~\ref{sec:ablations} are reported, not used to choose settings. Hypotheses H1--H5 were registered in \texttt{proposal.md} before the main runs. H1, H4 and H5 are tested with paired $t$-tests over the five seeds at $\alpha=\rhval{const/alpha:2}$ (no correction for multiple comparisons, so they are descriptive); H2 and H3 are comparisons of means over seeds. The dropout decomposition (Section~\ref{sec:ablations}) was added after seeing the main results and is exploratory.

\textbf{Compute and provenance of the numbers.} Apple-silicon CPU (shared machine), 2 threads, torch CPU. One run takes a few seconds, and the study stayed within the 45-minute compute budget fixed in \texttt{research.yaml}. The logged runs are: main \rhval{count/main/runs}, oracle \rhval{count/abl_oracle/ts-oracle-shifted-val/runs}, dropout decomposition \rhval{count/abl_dropdet/dropout-trained-mlp-deterministic/runs}, ensemble-size sweep \rhval{count/sweep_members/runs} and dropout-rate sweep \rhval{count/sweep_dropout/runs} (three conditions, seeds 0--2), and one sanity run. The registry holds \rhval{count/runs} rows, because the main-group runs of the MLP and of TS are listed a second time in the oracle group, and those of the MLP and of MC dropout in the decomposition group (\texttt{rh log -{}-from-run}: copies of logged rows, no new computation), so that systems of different groups can be paired by seed. Every mean in the text and tables is an aggregate recomputed from the registry when the paper is built, and every difference, relative change and paired $t$-test $p$-value is a statistic of \texttt{rh compare}; none is typed by hand.
